\documentclass[12pt,a4paper]{article}
\usepackage{ctex}
\usepackage{amsmath}
\usepackage{graphicx}
\usepackage{listings}
\usepackage{xcolor}
\usepackage{geometry}
\usepackage{booktabs}
\usepackage{float}

% 页面设置
\geometry{left=2.5cm,right=2.5cm,top=2.5cm,bottom=2.5cm}

% 代码高亮设置
\lstset{
    language=Verilog,
    basicstyle=\ttfamily\small,
    keywordstyle=\color{blue},
    commentstyle=\color{green!60!black},
    stringstyle=\color{red},
    numbers=left,
    numberstyle=\tiny\color{gray},
    stepnumber=1,
    numbersep=5pt,
    backgroundcolor=\color{white},
    showspaces=false,
    showstringspaces=false,
    frame=single,
    tabsize=4,
    breaklines=true,
    breakatwhitespace=false
}

\title{ALU实验功能改进报告}
\date{}

\begin{document}

\maketitle


% \section{实验目的}
% 1. 深入理解ALU的工作原理和结构组成
% 2. 掌握Verilog HDL设计数字电路的方法
% 3. 对比独热码与二进制编码在控制信号设计中的优缺点
% 4. 学会在原有设计基础上进行功能扩展和优化
% 5. 掌握FPGA开发流程，包括代码编写、仿真验证和上板调试

\section{实验原理}
\subsection{ALU基本原理}
算术逻辑单元(ALU)是中央处理器的核心组成部分，负责执行所有的算术运算和逻辑运算。本实验中的原始ALU设计支持12种基本操作，包括加法、减法、比较、移位和逻辑运算等。

\subsection{编码方式改进原理}
原始设计采用12位独热码作为控制信号，每个操作对应一位控制信号。这种方式的优点是解码简单，但缺点是控制信号位数过多，浪费硬件资源。

4位二进制编码可以表示16种不同的状态，完全可以覆盖原有的12种操作，并且还有4个冗余状态可以用于扩展新功能。本实验将控制信号从12位减少到4位，通过组合逻辑解码产生各个操作的使能信号。

\subsection{新增操作}
1. 有符号比较大于置位 \\
\indent2. 按位同或 \\
\indent3. 低位加载

% 实验原理图（注释掉以防编译错误）
% \begin{figure}[H]
%     \centering
%     \includegraphics[width=0.9\textwidth]{fig5_3.png}
%     \caption{改进后的ALU实验原理图（参考实验指导书图5.3）}
%     \label{fig:alu_schematic}
% \end{figure}

\section{实验内容与修改}
% 本次实验主要对三个文件进行了修改：alu.v（核心ALU模块）、alu_display.v（显示与输入模块）和约束文件(.xdc)。

\subsection{操作码编码表设计}
保留原有12种操作的功能不变，重新分配4位二进制操作码，并新增三种操作，编码表如下：

\begin{table}[H]
\centering
\caption{改进后的ALU操作码编码表}
\begin{tabular}{cccccc}
\toprule
4位二进制码 & AC22(bit3) & AC23(bit2) & AB6(bit1) & W6(bit0) & ALU操作 \\
\midrule
0000 & 0 & 0 & 0 & 0 & 无操作 \\
0001 & 0 & 0 & 0 & 1 & 加法 \\
0010 & 0 & 0 & 1 & 0 & 减法 \\
0011 & 0 & 0 & 1 & 1 & 有符号比较，小于置位 \\
0100 & 0 & 1 & 0 & 0 & 无符号比较，小于置位 \\
0101 & 0 & 1 & 0 & 1 & 按位与 \\
0110 & 0 & 1 & 1 & 0 & 按位或非 \\
0111 & 0 & 1 & 1 & 1 & 按位或 \\
1000 & 1 & 0 & 0 & 0 & 按位异或 \\
1001 & 1 & 0 & 0 & 1 & 逻辑左移 \\
1010 & 1 & 0 & 1 & 0 & 逻辑右移 \\
1011 & 1 & 0 & 1 & 1 & 算术右移 \\
1100 & 1 & 1 & 0 & 0 & 高位加载 \\
1101 & 1 & 1 & 0 & 1 & 有符号比较，大于置位(新增) \\
1110 & 1 & 1 & 1 & 0 & 按位同或(新增) \\
1111 & 1 & 1 & 1 & 1 & 低位加载(新增) \\
\bottomrule
\end{tabular}
\label{tab:opcode}
\end{table}

\begin{figure}[H]
    \centering
    \includegraphics[width=1\textwidth]{原理图.jpg}
    \caption{原理图}

\end{figure}

\subsection{alu.v模块修改}
这是本次实验的核心修改文件，主要修改内容如下：

1. 模块端口修改：将12位独热码控制信号改为4位二进制控制信号
\begin{lstlisting}
// 修改前
module alu(
    input  [11:0] alu_control,  // ALU控制信号
    input  [31:0] alu_src1,     // ALU操作数1,为补码
    input  [31:0] alu_src2,     // ALU操作数2，为补码
    output [31:0] alu_result    // ALU结果
    );

// 修改后
module alu(
    input  [3:0] alu_control,   // ALU控制信号(4位二进制)
    input  [31:0] alu_src1,     // ALU操作数1,为补码
    input  [31:0] alu_src2,     // ALU操作数2，为补码
    output [31:0] alu_result    // ALU结果
    );
\end{lstlisting}

2. 控制信号解码修改：移除独热码解码，添加新操作信号并实现4位二进制解码
\begin{lstlisting}
// 新增三个操作信号
wire alu_sgt;   //有符号比较，大于置位(新增)
wire alu_xnor;  //按位同或(新增)
wire alu_lli;   //低位加载(新增)

// 4位二进制控制信号解码
assign alu_add  = (alu_control == 4'b0001);
assign alu_sub  = (alu_control == 4'b0010);
assign alu_slt  = (alu_control == 4'b0011);
assign alu_sltu = (alu_control == 4'b0100);
assign alu_and  = (alu_control == 4'b0101);
assign alu_nor  = (alu_control == 4'b0110);
assign alu_or   = (alu_control == 4'b0111);
assign alu_xor  = (alu_control == 4'b1000);
assign alu_sll  = (alu_control == 4'b1001);
assign alu_srl  = (alu_control == 4'b1010);
assign alu_sra  = (alu_control == 4'b1011);
assign alu_lui  = (alu_control == 4'b1100);
assign alu_sgt  = (alu_control == 4'b1101); // 新增
assign alu_xnor = (alu_control == 4'b1110); // 新增
assign alu_lli  = (alu_control == 4'b1111); // 新增
\end{lstlisting}

3. 新增操作实现：添加三个新操作的结果计算逻辑
\begin{lstlisting}
// 按位同或结果为异或结果按位取反
assign xnor_result = ~xor_result;             
// 低位加载：立即数放在低半字节
assign lli_result = {16'd0, alu_src2[15:0]};  

// 有符号大于置位结果
assign sgt_result[31:1] = 31'd0;
assign sgt_result[0]    = (~alu_src1[31] & alu_src2[31]) | 
                          (~(alu_src1[31]^alu_src2[31]) & ~adder_result[31]);
\end{lstlisting}

4. 结果选择逻辑修改
\begin{lstlisting}
// 选择相应结果输出
    assign alu_result = 
        (alu_control == 4'b0001) ? add_sub_result :
        (alu_control == 4'b0010) ? add_sub_result :
        (alu_control == 4'b0011) ? slt_result :
        (alu_control == 4'b0100) ? sltu_result :
        (alu_control == 4'b0101) ? and_result :
        (alu_control == 4'b0110) ? nor_result :
        (alu_control == 4'b0111) ? or_result :
        (alu_control == 4'b1000) ? xor_result :
        (alu_control == 4'b1001) ? sll_result :
        (alu_control == 4'b1010) ? srl_result :
        (alu_control == 4'b1011) ? sra_result :
        (alu_control == 4'b1100) ? lui_result :
        (alu_control == 4'b1101) ? sgt_result :
        (alu_control == 4'b1110) ? xnor_result :
        (alu_control == 4'b1111) ? lli_result :
        32'd0;
\end{lstlisting}

\subsection{alu\_display.v模块修改}
1. 添加拨码开关输入端口：添加四个拨码开关作为ALU控制信号输入
\begin{lstlisting}
// 新增：4位拨码开关作为ALU控制信号(从高到低)
input alu_ctrl3, // AC22 - 最高位
input alu_ctrl2, // AC23
input alu_ctrl1, // AB6
input alu_ctrl0, // W6 - 最低位
\end{lstlisting}

2. 修改控制信号输入逻辑：优先使用拨码开关输入，触摸屏输入作为备用
\begin{lstlisting}
// ALU控制信号：优先使用拨码开关输入，触摸屏输入作为备用
always @(posedge clk)
begin
    if (!resetn)
    begin
        alu_control <= 4'd0;
    end
    else if (input_valid && input_sel==2'b00)
    begin
        alu_control <= input_value[3:0]; // 触摸屏输入控制信号
    end
    else
    begin
        alu_control <= {alu_ctrl3, alu_ctrl2, alu_ctrl1, alu_ctrl0}; // 拨码开关输入
    end
end
\end{lstlisting}

3. 修改控制信号显示：调整显示位宽以正确显示4位控制信号
\begin{lstlisting}
6'd3 :
begin
    display_valid <= 1'b1;
    display_name  <= "CONTR";
    display_value <={28'd0, alu_control}; // 显示4位控制信号
end
\end{lstlisting}

\subsection{约束文件修改}
在约束文件中添加四个拨码开关的引脚约束：
\begin{lstlisting}
# 新增：ALU控制信号拨码开关(从高到低)
set_property PACKAGE_PIN AC22 [get_ports alu_ctrl3]
set_property PACKAGE_PIN AC23 [get_ports alu_ctrl2]
set_property PACKAGE_PIN AB6  [get_ports alu_ctrl1]
set_property PACKAGE_PIN W6   [get_ports alu_ctrl0]

set_property IOSTANDARD LVCMOS33 [get_ports alu_ctrl3]
set_property IOSTANDARD LVCMOS33 [get_ports alu_ctrl2]
set_property IOSTANDARD LVCMOS33 [get_ports alu_ctrl1]
set_property IOSTANDARD LVCMOS33 [get_ports alu_ctrl0]
\end{lstlisting}

% \section{波形仿真验证}
% 为了验证改进后的ALU功能正确性，我们使用ModelSim进行了波形仿真。

% \subsection{仿真测试平台设计}
% 设计了一个简单的测试平台，依次设置不同的操作码和操作数，观察ALU的输出结果是否正确。测试用例覆盖了所有15种操作，重点验证了新增的三种操作。

% \subsection{仿真结果分析}
% % 仿真波形图（注释掉以防编译错误）
% % \begin{figure}[H]
% %     \centering
% %     \includegraphics[width=1\textwidth]{alu_simulation.png}
% %     \caption{ALU功能仿真波形图}
% %     \label{fig:simulation}
% % \end{figure}

% 仿真结果表明：
% 1. 原有12种操作功能完全正确，与改进前的结果一致
% 2. 新增的有符号大于置位操作：当alu\_src1=5，alu\_src2=3时，结果为0x00000001；当alu\_src1=3，alu\_src2=5时，结果为0x00000000
% 3. 新增的按位同或操作：当alu\_src1=0xAAAA5555，alu\_src2=0x5555AAAA时，结果为0x00000000
% 4. 新增的低位加载操作：当alu\_src2=0x12345678时，结果为0x00005678
% 5. 所有操作的延迟均在一个时钟周期内，满足时序要求

\section{上箱验证}
\subsection{硬件连接}
按照约束文件的设置，将四个拨码开关AC22、AC23、AB6、W6分别连接到ALU控制信号的最高位到最低位。

\subsection{验证步骤}
1. 将编译好的比特流文件下载到FPGA开发板\\
\indent2. 通过拨码开关input\_sel选择输入源操作数1和源操作数2，使用触摸屏输入具体数值\\
\indent3. 通过四个控制拨码开关设置不同的操作码\\
\indent4. 观察触摸屏上显示的运算结果是否正确

\subsection{验证结果}
上箱验证结果与仿真结果完全一致：\\
\indent1. 所有原有操作均能正确执行\\
\indent2. 新增的三种操作功能正常\\
\indent3. 拨码开关控制灵敏，操作便捷\\
\indent4. 触摸屏显示清晰，输入输出正常

\begin{figure}[H]
    \centering
    \includegraphics[width=1\textwidth]{有符号比较大于置位.jpg}
    \caption{有符号比较大于置位功能}

\end{figure}
\begin{figure}[H]
    \centering
    \includegraphics[width=1\textwidth]{按位同或.jpg}
    \caption{按位同或功能}
\end{figure}

\begin{figure}[H]
    \centering
    \includegraphics[width=1\textwidth]{低位加载.jpg}
    \caption{低位加载功能}
\end{figure}



\end{document}